\chapter{CƠ SỞ LÝ THUYẾT}
\label{chap:coso}

\section{Tối ưu đa mục tiêu và quan hệ trội Pareto}

Để so sánh lời giải theo cùng một chiều, nhóm đưa ba mục tiêu về dạng cùng cực tiểu hoá:
\begin{equation}
  \mathbf{g}(\pi) = \big(g_1, g_2, g_3\big) = \big(f_1(\pi),\; f_2(\pi),\; 1 - f_3(\pi)\big).
\end{equation}
Lời giải $a$ \emph{trội} lời giải $b$, ký hiệu $a \prec b$, khi $a$ không tệ hơn $b$ ở mọi mục tiêu và
tốt hơn hẳn ở ít nhất một mục tiêu:
\begin{equation}
  a \prec b \iff \forall k:\ g_k(a) \le g_k(b) \;\;\text{và}\;\; \exists j:\ g_j(a) < g_j(b).
\end{equation}
Tập các lời giải không bị lời giải nào trội gọi là tập Pareto; ảnh của nó trong không gian mục tiêu
gọi là mặt Pareto (front). Mục đích của thuật toán là tìm một tập xấp xỉ mặt Pareto vừa gần mặt thật,
vừa trải đều.

Trong phần tìm kiếm cục bộ, phép so sánh dùng thêm một sai số $\varepsilon \ge 0$ để tránh chấp nhận
các thay đổi chỉ do sai số làm tròn số thực:
\begin{equation}
  a \prec_\varepsilon b \iff \forall k:\ g_k(a) \le g_k(b) + \varepsilon \;\;\text{và}\;\;
  \exists j:\ g_j(a) < g_j(b) - \varepsilon .
  \label{eq:dom-eps}
\end{equation}

\section{Thuật toán NSGA-II}

NSGA-II do Deb và cộng sự đề xuất năm 2002~\cite{deb2002}. Ý tưởng chính gồm ba phần.

\paragraph{Sắp xếp không trội nhanh.} Quần thể được chia thành các front $\mathcal{F}_1,
\mathcal{F}_2,\dots$: $\mathcal{F}_1$ là các cá thể không bị ai trội; bỏ $\mathcal{F}_1$ đi thì các cá
thể không bị ai trội còn lại tạo thành $\mathcal{F}_2$, cứ thế tiếp tục. Hạng của cá thể là chỉ số
front chứa nó. Với $M$ mục tiêu và $N$ cá thể, độ phức tạp là $O(MN^2)$.

\paragraph{Khoảng cách đông đúc.} Trong cùng một front, cá thể nằm ở vùng thưa được ưu tiên để giữ đa
dạng. Với mỗi mục tiêu $k$, sắp các cá thể theo $g_k$; hai cá thể ở hai đầu nhận khoảng cách vô cùng,
cá thể ở giữa được cộng thêm
\begin{equation}
  \frac{g_k(\text{sau}) - g_k(\text{trước})}{g_k^{\max} - g_k^{\min}} .
\end{equation}

\paragraph{Chọn lọc tinh hoa.} Cha mẹ và con được gộp thành $2N_{pop}$ cá thể, rồi lấy lần lượt từng
front cho đến khi đủ $N_{pop}$; front cuối cùng bị cắt theo khoảng cách đông đúc giảm dần. Cha mẹ được
chọn bằng giải đấu nhị phân: hạng thấp hơn thắng, cùng hạng thì khoảng cách lớn hơn thắng.

\paragraph{Trội có ràng buộc.} Nhóm xử lý giới hạn số xe bằng cách gán cho mỗi cá thể mức vi phạm
$CV = \max(0, m - K)$. Cá thể có $CV$ nhỏ hơn luôn trội; cùng $CV$ thì so theo trội Pareto. Cách này
không cần chọn hệ số phạt.

\section{Giải thuật lai và vai trò của khởi tạo, tìm kiếm cục bộ}

Thuật toán di truyền tìm kiếm tốt trên diện rộng nhưng tinh chỉnh chậm: lai ghép và đột biến là phép
biến đổi mù, không biết thay đổi nào làm lời giải tốt hơn. Hai cách bổ sung kinh điển là:
\begin{itemize}
  \item \textbf{Khởi tạo có định hướng}: đưa vào quần thể ban đầu một số lời giải dựng bằng heuristic,
        để thuật toán xuất phát gần vùng tốt thay vì từ các hoán vị ngẫu nhiên.
  \item \textbf{Tìm kiếm cục bộ sau biến đổi di truyền}: sau khi sinh con, cải thiện con bằng các bước
        nhỏ có kiểm tra mục tiêu. Thuật toán di truyền kết hợp tìm kiếm cục bộ thường gọi là giải thuật
        memetic~\cite{moscato1989}.
\end{itemize}
Cả hai đều có rủi ro. Khởi tạo quá thiên về heuristic làm mất đa dạng ban đầu. Tìm kiếm cục bộ tốn
thêm nhiều lần đánh giá, nên khi so với thuật toán gốc phải so cả theo cùng số lần đánh giá, không chỉ
theo cùng số thế hệ.

\section{Các heuristic khởi tạo cho bài toán định tuyến}

\paragraph{Láng giềng gần nhất (NN).} Xuất phát từ kho, mỗi bước đi tới khách chưa ghé gần điểm hiện tại
nhất. Thứ tự này cho tổng quãng đường ngắn ở phần đầu, có lợi cho $f_2$ và $F_3$, nhưng không xét khung
thời gian nên có thể đến muộn nhiều khách.

\paragraph{Sắp theo khung thời gian (TW).} Xếp khách theo giờ mở khung kỳ vọng $ETa$ tăng dần. Khách cần
nhận hàng sớm được phục vụ trước, có lợi cho $f_3$ và phần phạt $F_4$, nhưng có thể đi vòng xa.

\paragraph{Ngẫu nhiên.} Giữ phần lớn quần thể là hoán vị ngẫu nhiên như thuật toán gốc để duy trì đa dạng.

\section{Các phép tìm kiếm cục bộ trong một tuyến}

\paragraph{2-opt~\cite{croes1958}.} Chọn hai vị trí $i<j$ trong một tuyến và đảo ngược đoạn
$[i, j]$. Trên đường đi, phép này bỏ hai cạnh và nối lại theo cách khác, thường dùng để gỡ các đoạn đi
cắt chéo nhau.

\paragraph{Or-opt~\cite{or1976}.} Cắt một khối $L \in \{1,2,3\}$ khách liên tiếp, giữ nguyên thứ tự bên
trong khối, rồi chèn khối vào vị trí khác của cùng tuyến. Phép này bỏ ba cạnh và nối lại.

Cả hai phép đều không đổi tập khách của tuyến, chỉ đổi thứ tự ghé. Với bài toán có khung thời gian, đổi
thứ tự làm đổi giờ đến của các khách phía sau, nên ảnh hưởng đồng thời tới $F_1$ (tải trên từng chặng),
$F_3$ và $f_2$ (quãng đường), $F_4$ và $f_3$ (giờ đến).

\section{Tiêu chí chấp nhận bước đi khi có nhiều mục tiêu}

Tìm kiếm cục bộ một mục tiêu chỉ cần kiểm tra giá trị mới nhỏ hơn giá trị cũ. Với ba mục tiêu có hai
cách thông dụng.

\paragraph{Chấp nhận theo trội Pareto.} Chỉ nhận bước đi khi lời giải mới trội lời giải cũ
(công thức~\ref{eq:dom-eps}). Cách này an toàn: không bao giờ làm xấu mục tiêu nào. Nhược điểm là rất
khó thoả, vì một bước đi thường cải thiện mục tiêu này nhưng làm xấu nhẹ mục tiêu khác.

\paragraph{Chấp nhận theo tổng có trọng số.} Chọn vectơ trọng số $\mathbf{w} \ge 0$,
$\sum_k w_k = 1$ và vectơ hệ số chuẩn hoá $\mathbf{s} > 0$, rồi nhận bước đi khi
\begin{equation}
  S(\mathbf{g}) = \sum_{k=1}^{3} w_k\,\frac{g_k}{s_k}
  \quad\text{giảm hơn } \varepsilon .
  \label{eq:score}
\end{equation}
Hệ số chuẩn hoá cần thiết vì $f_1$ cỡ vài nghìn trong khi $1-f_3$ nhỏ hơn 1; không chuẩn hoá thì $f_1$
lấn át hoàn toàn. Mỗi vectơ trọng số hướng tìm kiếm về một vùng của mặt Pareto, nên bốc trọng số ngẫu
nhiên cho mỗi lần chạy tìm kiếm cục bộ giúp các con được đẩy về nhiều hướng khác nhau. Cách này dễ thoả
hơn, nhưng một bước được nhận không nhất thiết là cải thiện Pareto: một mục tiêu riêng lẻ có thể xấu đi.

\section{Các chỉ số đánh giá tập nghiệm}

Nhóm dùng các chỉ số sau, tính trên mục tiêu đã chuẩn hoá (\code{src/metrics.py}).
\begin{itemize}
  \item \textbf{HV}: thể tích vùng bị front trội, giới hạn bởi điểm tham chiếu $(1{,}1;\,1{,}1;\,1{,}1)$.
        Càng lớn càng tốt, tối đa $1{,}1^3 = 1{,}331$~\cite{zitzler1999}.
  \item \textbf{IGD}: khoảng cách trung bình từ mỗi điểm của tập tham chiếu tới điểm gần nhất của front.
        Càng nhỏ càng tốt.
  \item \textbf{SP}: độ lệch chuẩn của khoảng cách từ mỗi điểm tới điểm gần nó nhất. Càng nhỏ thì front
        càng đều.
\end{itemize}
Tập tham chiếu của một instance là phần không bị trội của hợp tất cả front cuối, gộp mọi cấu hình và
mọi seed. HV và IGD được tính bằng thư viện pymoo~\cite{pymoo}. Mỗi mục tiêu được chuẩn hoá về $[0,1]$ theo giá trị nhỏ nhất và lớn nhất trên mọi điểm của
mọi front. Vì tập tham chiếu và cận chuẩn hoá phụ thuộc toàn bộ các cấu hình được so, số HV chỉ có ý
nghĩa khi so giữa các cấu hình trong cùng một thí nghiệm và cùng một instance.

Để kiểm định khác biệt, nhóm dùng kiểm định Wilcoxon rank-sum trên 10 giá trị HV của hai cấu hình.
